\documentclass[11pt]{article}
\usepackage[margin=1in]{geometry}
\usepackage{amsmath,amssymb,amsthm}
\usepackage{booktabs}
\usepackage{enumitem}
\usepackage[hidelinks]{hyperref}
\setlist{itemsep=2pt,topsep=4pt}
\newtheorem{lemma}{Lemma}
\newtheorem{claim}{Claim}
\theoremstyle{definition}
\newtheorem{remark}{Remark}
\newcommand{\sig}{\sigma}
\newcommand{\E}{\mathbb{E}}
\renewcommand{\Pr}{\mathbb{P}}

\title{Generative model and core theory for \texttt{mcg\_popgen}\\[2pt]
\large Sequence $\to$ TF occupancy $\to$ methylation $\to$ (mutation, fitness) in a Wright--Fisher population}
\author{working notes}
\date{2026-10-01}

\begin{document}
\maketitle

\noindent This document states exactly what the simulations compute. Every quantity here has a counterpart in
\texttt{scripts/slim/*.slim} (SLiM~5.2) and \texttt{scripts/python/model.py}. Section~\ref{sec:cycle} gives the
per-generation algorithm, Sections~\ref{sec:mut}--\ref{sec:fit} the maps, Section~\ref{sec:models} the model set,
Section~\ref{sec:theory} the results used to validate the simulator (experiment E0), and Section~\ref{sec:e1} the
predictions for experiment E1.

\section{What the experiments are}
\begin{description}
\item[E0 (calibration).] No methylation or TF. It checks that the SLiM machinery (mutation matrix, rejection gating,
  selection, rescaling) reproduces quantities that can be derived exactly: neutral diversity, the site-frequency
  spectrum, Kimura fixation probabilities, and CpG decay. It says nothing biological. It certifies the simulator.
\item[E1 (passenger footprint).] Methylation is computed from sequence and can gate CpG mutation, but it
  \emph{never} enters fitness. Question: how much CpG enrichment near a selected TF motif arises with zero selection
  on any CpG?
\item[E2 (causal vs passenger; next).] Same maps; fitness reads methylation (model C) or occupancy (model P).
  Question: which observable statistics separate them once mean methylation is matched?
\end{description}

\section{State space and notation}
\begin{itemize}
\item A locus of $L$ bases (E1: $L=300$). A haplotype is $x=(x_0,\dots,x_{L-1})\in\{A,C,G,T\}^L$, coded $0,1,2,3$.
\item The population has $N$ diploid hermaphrodites, i.e.\ $2N$ haplotypes. Individual $j$ is a pair
  $(x^{(j,1)},x^{(j,2)})$. There is no recombination within the locus.
\item CpG set: $\mathcal{C}(x)=\{i: x_i=C,\ x_{i+1}=G\}$, indexed by the position $i$ of the C.
\item TF motif: a fixed 8-mer $\mu=\texttt{TGACTCAT}$ (it contains no CpG), at positions $M=\{s,\dots,s+7\}$, with
  $s=L/2-4$ and centre $c=s+3.5$.
\item Distance of position $i$ from the motif edge: $d_i=\max(0,\ |i-c|-4)$.
\item $\sig(z)=1/(1+e^{-z})$.
\end{itemize}
Methylation is \textbf{not} a state variable. It is recomputed from each haplotype's sequence whenever it is needed.
This encodes germline reprogramming: there is no transgenerational epigenetic inheritance, so methylation is a
deterministic molecular phenotype of the haplotype (in cis), averaged over cells.

\section{One generation (SLiM's Wright--Fisher cycle)}\label{sec:cycle}
Let generation $t$ consist of $N$ individuals with fitnesses $w_1,\dots,w_N$ (computed at the end of generation $t-1$).
\begin{enumerate}
\item \textbf{Parent choice.} Each of the $N$ offspring draws two parents independently, each with probability
  $w_j/\sum_k w_k$. Incidental selfing (the same parent drawn twice, probability $\approx 1/N$) is allowed, as in the
  standard WF model.
\item \textbf{Gamete formation.} Each parent transmits one of its two haplotypes, chosen with probability $1/2$.
  There is no recombination.
\item \textbf{Mutation.} The transmitted haplotype $x$ receives new mutations by the context-dependent process of
  Section~\ref{sec:mut}, including the methylation gate (evaluated on $x$, the parental sequence).
\item \textbf{Bookkeeping and sampling.} \texttt{late()} events record statistics. Fitness $w$ of each new individual
  is then computed from its two haplotypes (Section~\ref{sec:fit}) and used in step 1 of the next generation.
\end{enumerate}
Selection is therefore fecundity/viability selection applied once per generation, and drift is multinomial sampling of
$2N$ haplotypes. The effective size is $N_e=N$ (diploid WF), which is the convention for every $2N s$ and $4N u$ below.

\section{Mutation process}\label{sec:mut}
\subsection{Context-only (``ctx'') process}
Each haplotype base $i$ mutates to base $b\neq x_i$ at rate (per generation)
\begin{equation}
q_i(b\mid x)=\begin{cases} r\alpha & \text{if } (x_i,x_{i+1},b)=(C,G,T) \text{ or } (x_{i-1},x_i,b)=(C,G,A)\\
\alpha & \text{otherwise.}\end{cases}
\label{eq:q}
\end{equation}
This is Jukes--Cantor (total $u=3\alpha$ per site) plus CpG transition hypermutability $r$. In SLiM it is a $64\times4$
trinucleotide matrix (rows are $(x_{i-1},x_i,x_{i+1})$, columns are $b$). Bases off either end count as $A$ (SLiM's
convention), so position $0$ is never the G of a CpG.

\subsection{Methylation-gated (``meth'') process}
Deamination of 5mC is what makes CpGs hypermutable, so only the \emph{methylated} fraction carries the excess rate:
\begin{equation}
q_i^{\text{meth}}(T\mid x)=\alpha\,[1+(r-1)\,m_i(x)] \quad\text{for } i\in\mathcal{C}(x),
\label{eq:qmeth}
\end{equation}
and likewise for the G (G$\to$A) of the same CpG, using the same $m_i$. All other rates are as in~\eqref{eq:q}. An
unmethylated CpG ($m=0$) has the ordinary transition rate $\alpha$, and a fully methylated one has $r\alpha$.

\textbf{Implementation by thinning.} SLiM proposes CpG transitions at the maximal rate $r\alpha$ (from the matrix). A
\texttt{mutation()} callback then accepts each proposal with probability
\begin{equation}
g(m)=\frac{1+(r-1)m}{r}\in[1/r,\,1].
\end{equation}
\begin{lemma}[thinning]
If events occur as a Poisson process of rate $\lambda$ and each is kept independently with probability $g$, the kept
events form a Poisson process of rate $g\lambda$. Hence accepted CpG transitions occur at exactly
rate~\eqref{eq:qmeth}.
\end{lemma}
SLiM already uses thinning internally: it proposes at the matrix's maximal row rate and rejects according to the actual
context. The gate is a second, state-dependent thinning. One caveat matters. SLiM passes the \emph{parental}
haplotype to the callback, so $m_i$ is evaluated on the sequence being copied. That is exact only without
recombination, which is why the methylation models have none. E0 verified that a constant gate $g=0.5$ at $r=20$ is
indistinguishable from ctx at $r=10$.

\section{Maps: sequence $\to$ occupancy $\to$ methylation}\label{sec:maps}
\textbf{TF occupancy} (per haplotype). With $S(x)=\sum_{k=0}^{7}\mathbf{1}[x_{s+k}=\mu_k]$ the number of matches,
\begin{equation}
O(x)=\sig\!\big(\beta\,(S(x)-S_0)\big),\qquad \beta=2,\ S_0=6.
\end{equation}
This gives $O=0.982,\,0.881,\,0.5,\,0.119$ for $S=8,7,6,5$. It is the simplest stand-in for a PWM/binding-energy
model: occupancy saturates for good sites and collapses after two mismatches.

\textbf{Methylation} of the CpG at $i\in\mathcal{C}(x)$:
\begin{equation}
m_i(x)=\sig\!\Big(a_0-a_1\,\rho_i(x)-a_2\,O(x)\,e^{-d_i/\lambda}\Big),\qquad
\rho_i(x)=\#\{j\in\mathcal{C}(x)\setminus\{i\}: |j-i|\le W\}.
\label{eq:meth}
\end{equation}
\begin{itemize}
\item $a_0=3$ is the default state: $m=\sig(3)=0.953$ (bulk mammalian CpGs are mostly methylated).
\item $a_2=6$ and $\lambda=25$ set TF protection, i.e.\ low-methylated regions around bound sites (Stadler et al.\ 2011).
  A CpG next to a bound motif has $m=\sig(3-6\cdot0.982)=0.053$.
\item $a_1$ and $W$ set CpG-density feedback: dense CpG clusters resist methylation (CGI-like). With $a_1=0$ the map is
  local and per-CpG. With $a_1>0$ it is \emph{degenerate}: many sequences give the same $m$, and selection on $m$
  spreads over many sites.
\end{itemize}
A Python mirror (\texttt{model.py}) recomputes $m_i$ from SLiM's output sequences. Agreement is $<10^{-6}$, the limit
of printed precision.

\section{Expression and fitness}\label{sec:fit}
Diploid expression is additive (codominant) across the two haplotypes, and fitness is stabilising around the
optimum $E^*=1$:
\begin{equation}
E_j=\tfrac12\big[g(x^{(j,1)})+g(x^{(j,2)})\big],\qquad w_j=\exp\!\big(-\kappa\,(1-E_j)^2\big).
\end{equation}
The haplotype readout $g(\cdot)$ is what distinguishes the models:
\begin{equation}
\text{P (passenger): } g(x)=O(x);\qquad \text{C (causal, E2): } g(x)=1-\bar m_{\text{prom}}(x);\qquad
\text{M0 (neutral): } \kappa=0.
\end{equation}
\textbf{Optimum.} E1 uses $E^*=1$. Because $O\le0.982$, that is directional selection with diminishing returns,
not true stabilising selection. E2 uses an intermediate optimum (e.g.\ $E^*=0.5$). Then many genotypes sit at the
optimum, which allows compensatory turnover and lets methylation be conserved while sequence diverges.

\textbf{Reporting selection.} $\kappa$ by itself is not interpretable. For a genotype $G$ we report
$2N s_G$, with $s_G=1-w_G/w_{8/8}$. $\kappa$ is set so that one heterozygous motif mismatch has $2Ns_{8/7}=10$. That gives
$\kappa=11.7$ at $N=100$ and $2Ns=30,\ 108,\ 181$ for $7/7$, $8/6$, $8/5$. The motif is under strong selection and
CpGs are under none.

\section{The model set and the causal/passenger question}\label{sec:models}
\begin{center}
\begin{tabular}{@{}lllll@{}}
\toprule
model & fitness reads & mutation & CpG feedback & role\\
\midrule
M0\_ctx & nothing & ctx & -- & neutral baseline (E0's process)\\
M0\_meth & nothing & meth & $a_1=0$ & motif decays, so no protection\\
M0\_meth\_fb & nothing & meth & $a_1=0.5$ & can CpG islands self-sustain?\\
P\_ctx & $O$ & ctx & -- & selection on TF, no gating\\
P\_meth & $O$ & meth & $a_1=0$ & \textbf{passenger footprint}\\
P\_meth\_fb & $O$ & meth & $a_1=0.5$ & passenger + feedback\\
C (E2) & $\bar m$ & meth & swept & methylation causal\\
\bottomrule
\end{tabular}
\end{center}
\begin{remark}[why P vs C is a statement about the genotype--fitness map]
Because $m$ is a function of $x$, every model is ``selection on sequence''. P and C share both maps
($x\mapsto O$, $x\mapsto m$) and differ only in which composition, $w\circ O$ or $w\circ m$, defines fitness. A
mutation $x\to x'$ can discriminate them only if it changes one readout but not the other. Examples: a CpG gain or loss
outside the motif changes $m$ but not $O$; a motif mismatch with buffered methylation changes $O$ but barely changes
$m$. Such a mutation has a detectable fitness effect under C only if its selection is not swamped by drift, i.e.\
$2N\,|s(x\to x')|\gtrsim 1$. For a per-CpG effect $\Delta m$, $s\approx 2\kappa(1-E)\,\Delta m/2$ (to first order),
so identifiability requires
\begin{equation}
2N\kappa\,(1-E)\,|\Delta \bar m| \gtrsim 1 .
\end{equation}
Two consequences follow. Under a degenerate map ($a_1>0$, near saturation), single-CpG $|\Delta\bar m|$ is small,
so individual sites look neutral even when $\bar m$ is strongly conserved. And methylation-gated mutation (E1) can
create CpG patterns that look like constraint with no selection at all.
\end{remark}

\section{Theory used for validation (E0)}\label{sec:theory}
\subsection{Fixation probability}
With genotype fitnesses $1,1+s/2,1+s$, the diffusion has drift $M(p)=\tfrac{s}{2}p(1-p)$ and variance
$V(p)=p(1-p)/2N$. The scale function gives $u(p)=\dfrac{1-e^{-2Nsp}}{1-e^{-2Ns}}$, so a single new copy
($p=1/2N$) fixes with
\begin{equation}
u=\frac{1-e^{-s}}{1-e^{-2Ns}}\quad\xrightarrow{s\to0}\quad\frac{1}{2N}.
\end{equation}
E0: 20{,}000 trials at each of 8 settings, all within 1.3 SE.

\subsection{Neutral diversity at finite sites}
Two lineages coalesce after $T\sim\text{Exp}(1/2N)$ generations, so they are separated by total branch length $2T$.
Under JC, the probability that a site is identical after branch length $\ell$ is $\tfrac14+\tfrac34e^{-4u\ell/3}$.
Then
\begin{equation}
\pi=1-\E\!\left[\tfrac14+\tfrac34e^{-8uT/3}\right]=\frac{\theta}{1+4\theta/3},\qquad \theta=4Nu.
\end{equation}
With $\theta=0.01$, $\pi=0.009868$. E0 matched this at $N=500,250,100$.

\subsection{Neutral SFS}
In a sample of $n$, $\E[\xi_i]=\theta/i$. Folding gives the shape $\propto 1/i+1/(n-i)$. E0 obtained $\chi^2$ of 24--33
on 24 degrees of freedom.

\subsection{The single-lineage oracle}
\begin{lemma}
Assume no selection and a mutation process in which each haplotype's rates depend only on its own sequence (true for
ctx and meth, since methylation is cis). Then, for any haplotype sampled at generation $t$,
$\Pr(x(t)=y)=\big[e^{tQ}\big]_{x(0),y}$, where $Q$ is the generator of the per-sequence mutation chain on
$\{A,C,G,T\}^L$.
\end{lemma}
\begin{proof}[Sketch]
Trace the sampled haplotype's ancestry back to generation $0$. Under neutrality, the genealogy (which parent and which
haplotype was copied) is independent of the sequences. Along the ancestral path, the sequence changes only by
mutation, with rates that depend only on the current sequence on that path. So $x(t)$ is the chain run for $t$ steps
from $x(0)$. Discrete versus continuous time differs at $O(\text{rate}^2)$.
\end{proof}
Hence $\E[\#\text{CpG}](t)$ in the population equals that of a Gillespie simulation of one sequence. This is the
\texttt{oracle.py} check, and it needs no equilibrium. Under selection the lemma fails, because genealogy and
sequence become coupled. For P models the oracle instead freezes the motif, and SLiM minus oracle then measures the
effect of motif polymorphism.

\subsection{Rescaling}
The diffusion limit depends on parameters only through $Nu$, $Ns$, $Nr_{\text{rec}}$ and time in units of $N$. Running
$(N/Q,\ uQ,\ sQ,\ t/Q)$ preserves these and costs $Q^2$ less. This is valid while per-generation rates stay small
($r\alpha Q\ll1$, $sQ\ll1$). E0 checked neutral statistics at $Q=1,2,5$. E1 checks P\_meth at $Q=1,2,4$.

\section{E1 prediction: a two-state approximation}\label{sec:e1}
Treat one dinucleotide position as either CpG or not. A CpG with methylation $m$ is lost at rate
\begin{equation}
\ell(m)=\underbrace{2\alpha[1+(r-1)m]}_{\text{C}\to\text{T},\ \text{G}\to\text{A}}+\underbrace{4\alpha}_{\text{other changes of C or G}}.
\end{equation}
A non-CpG becomes a CpG at rate $\gamma\alpha$, where $\gamma\approx 6/15$: 6 of the 15 non-CpG dinucleotides are one
step from CG (CA, CC, CT, AG, GG, TG), if dinucleotides were uniform. The equilibrium CpG frequency per position is then
\begin{equation}
f(m)\approx\frac{\gamma\alpha}{\gamma\alpha+\ell(m)}.
\end{equation}
Checks and predictions at $r=20$:
\begin{itemize}
\item $m=1$ (E0's ctx): $f\approx0.4/44.4=0.009$, i.e.\ 0.90 CpG per 100\,bp. E0 measured 0.94. The approximation is
  good because the TG/CA decay products are themselves one step from CG.
\item Far flank in P\_meth ($m=0.953$): $\ell=42.2\alpha$, giving 0.94 per 100\,bp.
\item Next to the bound motif ($m=0.053$): $\ell=8.0\alpha$, giving 4.8 per 100\,bp, a \textbf{${\sim}5\times$ enrichment
  with zero selection on CpGs}. It falls to ${\sim}1.3\times$ at $d=25$, where $m=0.70$.
\end{itemize}
This ignores dependence between neighbouring positions and the motif's own polymorphism. The oracle computes the exact
neutral-flank expectation, and SLiM includes everything.

\paragraph{The feedback case.} With $a_1>0$, retention lowers $m$ (through $\rho$), which raises retention further. The
map $f\mapsto f(m(\rho(f)))$ can have two stable fixed points: a methylated, CpG-poor state and an unmethylated,
CpG-rich island. Whether an island self-sustains after the TF is lost (M0\_meth\_fb) or needs the TF to seed it
(P\_meth\_fb) is an empirical question for E1.

\section{P0: per-CpG priors across species}\label{sec:p0}
Extend the two-state chain of Section~\ref{sec:e1} with selection $S=2Ns$ favouring the CpG (from selection on
the CpG, or from gBGC). Under origin-fixation, the loss and gain rates become
\begin{equation}
a=\alpha\,\ell(m)\,\varphi(-S),\qquad b=\alpha\,\gamma\,\varphi(S),\qquad \varphi(S)=\frac{S}{1-e^{-S}},
\end{equation}
where $\varphi$ is $2N\times$ Kimura's $u$ relative to neutral, and $m$ is \emph{germline} methylation. Because
$\varphi(S)/\varphi(-S)=e^{S}$, the stationary CpG frequency is $f=\gamma e^{S}/(\gamma e^{S}+\ell(m))$: the neutral
stationary distribution weighted by $e^{S\cdot\text{state}}$, as in Sella \& Hirsh. For two species separated by
total branch length $T$, measured through neutral non-CpG divergence $D=3\alpha T$, reversibility gives
\begin{equation}
\Pr(\text{CpG in B}\mid\text{CpG in A})=f+(1-f)\,e^{-(a+b)T}.
\end{equation}
\textbf{Polymorphism.} For new mutations with $\sigma=2Ns$, the Poisson-random-field (Sawyer--Hartl) density of
derived-allele frequency, per unit $\theta$, is
\begin{equation}
\phi_\sigma(x)=\frac{1-e^{-\sigma(1-x)}}{(1-e^{-\sigma})\,x(1-x)}\quad\xrightarrow{\sigma\to0}\quad\frac1x .
\end{equation}
Expected segregating sites in a sample of $n$ are $\theta\int_0^1\phi_\sigma(x)\,[1-x^n-(1-x)^n]\,dx$; SLiM agrees
at $\sigma=0,-1,-2,-4,-8$. A change in germline $m$ rescales polymorphism and divergence of CpG-loss alleles by the
same factor $\ell(m)$, so $P/D$ is unchanged. Selection gives
$P/D \propto \int\phi_{-S}/\varphi(-S)$, which is 1.5, 2.4 and 7.8 at $S=1,2,4$. This is the McDonald--Kreitman-type
contrast that separates mutation from selection, where divergence and o/e alone lie on a ridge in $(m,S)$.

\section{Assumptions made explicit}
\begin{enumerate}
\item Methylation is regenerated from sequence each generation. There are no epialleles.
\item Methylation is a cell-averaged deterministic value. Stochastic cell-to-cell noise would matter only through
  nonlinearity in fitness.
\item Cis only. There are no trans factors, so the TF's concentration is constant.
\item There is no methylation $\to$ TF arrow (methyl-sensitive TFs such as NRF1 are deferred). This is what makes P a
  ``pure'' passenger.
\item No recombination and no gBGC within the locus.
\item Constant $N$, a single population, and no demography. Two-population splits come in E3.
\item Codominant, additive expression. Fitness is Gaussian in $E$ with the optimum at the maximum.
\end{enumerate}

\section{Parameters (E1 baseline, $Q=2$)}
\begin{center}
\begin{tabular}{@{}lll@{}}
\toprule
symbol & value & meaning\\
\midrule
$N$ & 100 & diploids ($Q=1$: 200)\\
$L$ & 300 & bp\\
$\alpha$ & $2\times10^{-5}$ & per target per generation; $u=3\alpha$, $\theta=4Nu=0.024$\\
$r$ & 20 (5, 50) & CpG transition multiplier\\
$\beta,S_0$ & 2, 6 & occupancy curve\\
$a_0,a_1,a_2$ & 3, 0 or 0.5, 6 & methylation map\\
$\lambda,W$ & 25 (10, 50), 25 & protection range, density window (bp)\\
$\kappa$ & 11.7 & so that $2Ns_{8/7}=10$\\
burn-in, measurement & $5/u$, $2/u$ generations & composition equilibrates on the $1/u$ timescale, not $N$\\
\bottomrule
\end{tabular}
\end{center}
The burn-in is set by $1/u$ rather than by $N$. Base composition (CpG content) changes only by substitution, which
takes $\sim1/u$ generations per site regardless of $N$. A $10N$ burn-in equilibrates diversity but not CpG content.

\end{document}
